\chapter{Metodyka oceny modeli}

\section{Cel rozdziału}

Celem tego rozdziału jest określenie zasad, według których oceniane są systemy rankingowe, metamodel oraz benchmark bukmacherski. Jest to istotne, ponieważ badany problem nie polega wyłącznie na wskazaniu zwycięzcy meczu, lecz na oszacowaniu prawdopodobieństwa zwycięstwa jednej ze stron. Z tego względu metryki powinny oceniać jakość predykcji probabilistycznej, a nie tylko decyzję klasyfikacyjną po przyjęciu arbitralnego progu.

W pracy przyjęto, że każdy model zwraca wartość \(\hat{p}_i \in (0,1)\), interpretowaną jako prawdopodobieństwo zwycięstwa pierwszej strony meczu w obserwacji \(i\). Rzeczywisty wynik oznaczono jako \(y_i \in \{0,1\}\), gdzie \(1\) oznacza zwycięstwo pierwszej strony, a \(0\) zwycięstwo drugiej strony.

\section{Wspólna próba ewaluacyjna}

Główna ewaluacja modeli w pracy jest prowadzona na wspólnej próbie meczów, dla których dostępne są zarówno dane sportowe z GOL.GG, jak i historyczne kursy bukmacherskie pochodzące z OddsPortal. Takie ograniczenie próby pozwala porównywać modele sportowe z benchmarkiem rynkowym w tych samych warunkach. Dzięki temu różnice w metrykach nie wynikają z odmiennego zakresu danych, lecz z jakości analizowanych predykcji.

Pełniejsza historia sportowa może być używana do budowy rankingów i historycznych cech kontekstowych, ponieważ w realistycznym scenariuszu predykcyjnym wcześniejsze wyniki zespołów i zawodników są znane przed meczem. Końcowe porównanie wyników odbywa się jednak na próbie wspólnej z rynkiem bukmacherskim. Jest to szczególnie ważne dla hipotezy dotyczącej konkurencyjności modelu względem kursów otwarcia i zamknięcia.

\section{Zasada chronologicznej walidacji}

Wszystkie modele oceniane w pracy muszą respektować kolejność czasową obserwacji. Oznacza to, że predykcja dla danego meczu może wykorzystywać wyłącznie informacje dostępne przed jego rozpoczęciem. Po zakończeniu meczu jego wynik może zostać użyty do aktualizacji rankingów lub cech historycznych, ale wyłącznie dla kolejnych obserwacji.

Takie podejście ogranicza ryzyko przecieku informacji. W szczególności niedopuszczalne byłoby używanie statystyk z aktualnie rozgrywanej gry, takich jak przewaga w złocie w piętnastej minucie, do predykcji wyniku przedmeczowego. Zmiennymi dopuszczalnymi są natomiast agregaty historyczne, obliczane na podstawie wcześniejszych meczów.

\section{Bootstrap blokowy dla porównania modeli}

Oprócz punktowych wartości metryk w pracy wykorzystano bootstrap blokowy do oceny stabilności różnic pomiędzy modelami. Jednostką losowania był miesiąc, ponieważ mecze z tego samego okresu mogą dotyczyć podobnych wersji gry, turniejów, składów i warunków rynkowych. Zwykły bootstrap losujący pojedyncze mecze mógłby zatem zaniżać niepewność przez ignorowanie zależności czasowych.

Dla każdego meczu liczono różnicę strat LogLoss pomiędzy modelem referencyjnym a modelem porównywanym, a następnie losowano miesięczne bloki z powtórzeniami i obliczano średnią różnicę w każdej próbie bootstrapowej. W ten sposób wyznaczano empiryczny przedział ufności dla różnicy jakości modeli. Znak różnicy definiowano osobno w każdym eksperymencie, tak aby interpretacja tabeli była jednoznaczna.

Bootstrap blokowy nie zastępuje walidacji chronologicznej, lecz ją uzupełnia. Walk-forward pokazuje wynik modelu w kolejnych okresach historycznych, natomiast bootstrap sprawdza, czy obserwowana różnica LogLoss jest stabilna względem zmian w składzie miesięcy ewaluacyjnych. Dlatego wybór finalnego modelu nie opiera się wyłącznie na pojedynczej wartości LogLoss, ale także na przedziałach ufności różnic względem wariantów alternatywnych i benchmarków.

Procedurę walk-forward stosowaną w eksperymentach można zapisać jako następującą sekwencję kroków:

\begin{enumerate}
    \item Uporządkuj mecze chronologicznie według daty rozegrania.
    \item Wybierz początkowy zbiór historyczny, który jest dostępny przed pierwszym okresem testowym.
    \item Wytrenuj model wyłącznie na danych historycznych dostępnych w danym momencie.
    \item Wygeneruj predykcje dla kolejnych $N$ meczów, bez używania ich wyników ani statystyk pomeczowych.
    \item Zapisz predykcje, rzeczywiste wyniki oraz wartości strat, np. LogLoss, do późniejszej analizy.
    \item Po zakończeniu tych $N$ meczów dopisz je do historii i zaktualizuj rankingi oraz cechy kroczące.
    \item Powtarzaj kroki 3--6 aż do końca okresu ewaluacyjnego.
    \item Po zebraniu wszystkich predykcji out-of-time pogrupuj błędy według miesięcy.
    \item Losuj miesięczne bloki z powtórzeniami, aby uzyskać bootstrapowy rozkład różnic LogLoss.
    \item Wyznacz przedział ufności dla $\Delta$ LogLoss i wykorzystaj go do oceny stabilności przewagi jednego modelu nad drugim.
\end{enumerate}

Taki zapis lepiej oddaje praktyczną implementację eksperymentu: model nie otrzymuje jednego losowego podziału danych, lecz przechodzi po kolejnych fragmentach historii. Każda predykcja jest więc predykcją out-of-time, a bootstrap blokowy jest wykonywany dopiero po zakończeniu całej procedury walk-forward.

\section{Metryki oceny}

\subsection{LogLoss}

Podstawową metryką wyboru modelu jest LogLoss, czyli logarytmiczna strata dla klasyfikacji binarnej. Dla zbioru \(N\) obserwacji definiuje się ją następująco:

\begin{equation}
    \operatorname{LogLoss} = -\frac{1}{N}\sum_{i=1}^{N}\left[y_i \log(\hat{p}_i) + (1-y_i)\log(1-\hat{p}_i)\right].
\end{equation}

Metryka ta silnie karze predykcje pewne, ale błędne. Jeżeli model przypisuje bardzo wysokie prawdopodobieństwo zdarzeniu, które nie następuje, strata jest duża. Jest to pożądane w tej pracy, ponieważ model ma zwracać wiarygodne prawdopodobieństwa, a nie tylko klasyfikować mecze jako zwycięstwo lub porażkę. LogLoss należy do właściwych reguł punktacji, czyli zachęca model do raportowania prawdopodobieństw zgodnych z rzeczywistym przekonaniem predykcyjnym \parencite{gneiting2007}.

\subsection{AUC}

AUC mierzy zdolność modelu do rozróżniania zwycięskich i przegranych obserwacji. Intuicyjnie można ją interpretować jako prawdopodobieństwo, że losowo wybrany wygrany przypadek otrzyma wyższy wynik modelu niż losowo wybrany przypadek przegrany. Wartość \(0{,}5\) odpowiada klasyfikatorowi losowemu, a wartość \(1\) oznacza idealne uporządkowanie przypadków.

AUC jest użyteczna do oceny siły sygnału sportowego, ponieważ pokazuje, czy model poprawnie porządkuje mecze według szans zwycięstwa. Nie jest jednak wystarczającą metryką wyboru modelu probabilistycznego. Model może mieć wysokie AUC, ale jednocześnie generować źle skalibrowane prawdopodobieństwa.

\subsection{Brier Score}

Brier Score mierzy średni kwadratowy błąd predykcji probabilistycznej:

\begin{equation}
    \operatorname{Brier} = \frac{1}{N}\sum_{i=1}^{N}(\hat{p}_i - y_i)^2.
\end{equation}

Niższa wartość oznacza lepszą jakość predykcji. W porównaniu z LogLoss metryka ta jest mniej surowa wobec bardzo pewnych błędów, ale jest łatwiejsza interpretacyjnie, ponieważ bezpośrednio mierzy odległość przewidywanego prawdopodobieństwa od rzeczywistego wyniku.

\subsection{ECE}

ECE, czyli expected calibration error, mierzy kalibrację prawdopodobieństw. Predykcje dzieli się na przedziały prawdopodobieństwa, a następnie porównuje średnie przewidywane prawdopodobieństwo z empirycznym odsetkiem zwycięstw w każdym przedziale. Ogólna postać metryki jest następująca:

\begin{equation}
    \operatorname{ECE} = \sum_{b=1}^{B}\frac{|S_b|}{N}\left|\operatorname{acc}(S_b) - \operatorname{conf}(S_b)\right|,
\end{equation}

gdzie \(S_b\) oznacza zbiór obserwacji w przedziale \(b\), \(\operatorname{acc}(S_b)\) empiryczny odsetek zwycięstw, a \(\operatorname{conf}(S_b)\) średnie przewidywane prawdopodobieństwo. Niska wartość ECE oznacza, że predykcje modelu są dobrze skalibrowane. Przykładowo, wśród meczów ocenionych przez model na około \(70\%\) rzeczywisty odsetek zwycięstw powinien być zbliżony do \(70\%\).

\subsection{Accuracy}

Accuracy, czyli odsetek poprawnych klasyfikacji, jest w tej pracy metryką wyłącznie pomocniczą. Aby ją obliczyć, należy przyjąć próg decyzyjny, najczęściej \(0{,}5\). Taki próg jest jednak arbitralny i nie wykorzystuje pełnej informacji zawartej w predykcji probabilistycznej. Dwa modele mogą mieć podobną accuracy, ale bardzo różną jakość prawdopodobieństw. Z tego powodu accuracy nie stanowi głównego kryterium wyboru modelu.

\section{Hierarchia wyboru modelu}

W pracy przyjęto hierarchię metryk dopasowaną do celu badania. Ponieważ model ma generować prawdopodobieństwa porównywalne z prawdopodobieństwami implikowanymi przez kursy bukmacherskie, najważniejsza jest jakość probabilistyczna. Pojedyncze wartości LogLoss, ECE, Brier Score i AUC traktowano jednak przede wszystkim jako informację diagnostyczną, wskazującą najbardziej obiecujące warianty modelu.

Stabilność wyboru potwierdzano dopiero miesięcznym bootstrapem blokowym różnic LogLoss. Jeżeli lepszy wynik punktowy nie był potwierdzony przedziałem ufności nieobejmującym zera, traktowano go jako sugestię, a nie jako stabilną poprawę. Hierarchia metryk określa więc, które własności modelu są najważniejsze, natomiast bootstrap blokowy sprawdza, czy obserwowana przewaga jest odporna na zmiany w składzie miesięcy ewaluacyjnych.

\begin{table}[H]
    \centering
    \caption{Hierarchia metryk wykorzystywanych przy wyborze modelu}
    \label{tab:metric_hierarchy}
    \begin{tabular}{p{0.18\textwidth}p{0.24\textwidth}p{0.46\textwidth}}
        \toprule
        Priorytet & Metryka & Rola w pracy \\
        \midrule
        1 & LogLoss & Główna metryka kierunkowa; wskazuje kandydatów do wyboru, a stabilność różnic jest następnie sprawdzana bootstrapem blokowym. \\
        2 & ECE & Kryterium pomocnicze przy zbliżonym LogLoss; ocenia kalibrację prawdopodobieństw. \\
        3 & Brier Score & Dodatkowa miara błędu probabilistycznego, łatwiejsza interpretacyjnie niż LogLoss. \\
        4 & AUC & Ocena zdolności rozróżniania silniejszej strony meczu; nie decyduje samodzielnie o wyborze modelu. \\
        5 & Accuracy & Metryka pomocnicza, raportowana wyłącznie informacyjnie. \\
        \bottomrule
    \end{tabular}
\end{table}

Jeżeli dwa modele uzyskują bardzo podobny LogLoss, preferowany jest model lepiej skalibrowany, czyli o niższym ECE, o ile nie pogarsza to stabilności głównej metryki. Jeżeli różnice w LogLoss i ECE są niewielkie, pomocniczo analizowany jest Brier Score oraz AUC. Taka kolejność jest spójna z celem pracy: ważniejsze jest poprawne oszacowanie prawdopodobieństwa niż samo wskazanie bardziej prawdopodobnego zwycięzcy. Ostateczna interpretacja nie opiera się jednak na samym rankingu punktowych metryk, lecz na połączeniu: wartości LogLoss, kierunku różnicy względem alternatywy oraz przedziału ufności z bootstrapu blokowego.

\section{Konsekwencje dla interpretacji wyników}

Przyjęta hierarchia metryk wpływa na sposób interpretacji dalszych rozdziałów. Model o najwyższym AUC nie musi być modelem preferowanym, jeżeli jego LogLoss lub kalibracja są gorsze. Analogicznie, wysoka accuracy nie jest wystarczającym argumentem za wyborem modelu, ponieważ w problemie probabilistycznym istotne jest, czy wartości \(0{,}60\), \(0{,}70\) lub \(0{,}80\) odpowiadają rzeczywistej częstości zwycięstw.

Z tego powodu w końcowym porównaniu z benchmarkiem bukmacherskim szczególne znaczenie mają LogLoss, Brier Score i ECE. AUC pełni rolę uzupełniającą: pokazuje, czy model potrafi skutecznie porządkować mecze od mniej do bardziej prawdopodobnych zwycięstw pierwszej strony. Dzięki temu możliwe jest oddzielenie dwóch aspektów jakości modelu: siły rozróżniania oraz jakości kalibracji prawdopodobieństw.
